In this Master's Thesis, we study the problem of shared factors in RSA public keys extracted from \textit{Certificate Transparency logs}. This vulnerability may appear when key generation is not performed with sufficient randomness, causing different RSA moduli to share a prime factor and therefore become efficiently factorable.
\\
\\
This work is based, on the one hand, on the thesis \textit{Finding shared RSA factors in the Certificate Transparency logs}, which already addressed this problem using a classical \textit{batch GCD} approach, and, on the other hand, on Pelofske's article, which proposes the \textit{Binary Tree Batch GCD} algorithm as a more efficient alternative. Building on these references, we implement this algorithm in order to study its practical usefulness in this context.
\\
\\
To this end, we have developed an implementation in C++17 using GMP and OpenMP, which has already been validated in an initial stage on synthetically generated data. This phase has allowed us to verify that the algorithm correctly detects shared factors when they exist and to observe its runtime behavior as the size of the input set increases.
\\
\\
Subsequently, we have applied the same implementation to an initial real-world sample of RSA-2048 moduli extracted from \textit{CT logs}. Although this test is still preliminary and does not yet constitute a large-scale analysis of all available logs, it provides a useful validation of the procedure on real data.
\\
\\
In a later stage of the work, we will extend this analysis to a much larger collection of \textit{CT logs}, with the aim of better evaluating the scalability of the method and its usefulness for large-scale cryptographic audits. Overall, this thesis combines the study of previous work, methodological development, experimental implementation, and applied analysis of a real cybersecurity problem.